\subsection*{28. Source Actively When the Target Competence Is Observable in Public}

Some work leaves useful traces: well-maintained code, technical articles, open educational materials, published research, thoughtful documentation, conference explanations, or a history of responsible contributions to specialist communities. These artifacts can support targeted outreach. Yet an artifact proves only what it actually shows. Open-source contribution may indicate engineering habits, not teaching skill. A polished educational video may demonstrate presentation ability, not diagnostic feedback. Publications may demonstrate research participation without revealing the candidate's individual contribution.

A good sourcing message refers to \textbf{a specific, legitimate work artifact}, explains why the job may be relevant, and offers a low-pressure conversation. Avoid collecting personal information beyond what is necessary or repeatedly contacting people who decline. A technically accomplished candidate may value an outreach message that shows genuine familiarity with the work more than a long generic recruitment pitch.

Public evidence is a discovery tool. Formal selection should still assess the relevant competencies using fair, comparable methods.

\subsection*{29. Build Sourcing Experiments with Useful Denominators}

Suppose an illustrative recruiting experiment reaches 60 suitable professionals through professional communities, receives 18 expressions of interest, screens 12 against basic requirements, interviews 6, and makes 1 accepted offer. The initial contact-to-accepted-offer rate is $1/60$, but it is not an intrinsic property of that channel. It reflects offer attractiveness, role terms, market conditions, selection standards, and timing.

Track 4 distinct stages: \textbf{eligible reach}, \textbf{genuine interest}, \textbf{assessed qualification}, and \textbf{accepted offer}. Include the time to move between them. A channel generating many low-fit applications may consume more internal time than a channel with fewer, better-aligned applicants. Conversely, a channel that produces 1 exceptional candidate slowly may be unsuitable when delivery starts next week.

At small sample sizes, avoid declaring a precise conversion benchmark. Record explanations for withdrawals and declines. A pattern of candidates withdrawing after hearing compensation may identify a role-market mismatch, not a sourcing failure. A pattern of qualified women or regional candidates dropping out after an inconvenient test format deserves investigation even if the remaining candidates score highly.

\subsection*{30. Internal Mobility and Development Are Genuine Alternatives to Outside Hiring}

A firm may already have someone with the foundational capabilities who needs additional training or a redesigned position. Internal selection has information advantages: performance is observed over time, coworkers know real behavior, and the employer can examine past results. It also carries risks: managers may hoard talent, internal politics may bias assessment, and people may be typecast according to their existing position.

The same job-analysis discipline should apply. A strong support associate may be a promising operations manager but should not be promoted solely because of tenure or loyalty. Determine whether the new work requires different judgment, supervision, or technical competence. Use evidence of those capabilities and offer development where gaps are trainable.

Early Bayes may not have internal employees, but the principle matters over time. A rigorous junior reviewer could later become a faculty trainer if they demonstrate instructional diagnosis, coaching ability, and quality-control judgment. That is an institution deliberately creating its future candidate pool, rather than repeatedly competing for expensive senior hires.

\subsection*{31. The Candidate Experience Predicts Something About the Employer}

Applicants observe punctuality, organization, respect, payment practices, consistency, and whether the role changes during conversation. These experiences provide evidence about how the company manages people. An applicant asked to complete an unpaid 10-hour exercise may reasonably wonder whether the company values labor. A candidate who receives contradictory compensation promises from different interviewers may infer weak internal coordination. A company that never responds after an assessment burns reputational capital.

Not every candidate needs individualized coaching, but every candidate deserves clear expectations, truthful information, reasonable response timelines, and dignified treatment. A structured process can improve the experience if it makes criteria transparent. It can harm the experience if it becomes a robotic obstacle course with no explanation of why tasks matter.

Applicant reactions are not merely aesthetic. Researchers have studied the perceived fairness of selection methods, including job relatedness, respectful treatment, and timely communication [R9]. Candidates have choices; an unnecessarily adversarial process may selectively drive away people capable of finding better alternatives.

\subsection*{32. Sourcing and Selection Should Be Designed as 1 System}

An organization that recruits from a global community but tests candidates only through spontaneous spoken English in 1 accent has an incoherent system if the job mostly requires written technical review. It invites a broad pool and then evaluates a narrower, unrelated attribute. Likewise, a role advertised as remote cannot quietly require on-site attendance for every interview unless that requirement is truly necessary and disclosed.

Map the journey from 1st contact to decision. For each stage, ask what evidence the employer needs, how much candidate effort it requires, whether the format is accessible, who makes the judgment, and what happens if the candidate advances or withdraws. Candidates should not repeatedly prove the same competency because evaluators failed to share notes.

The output is not a marketing funnel with people reduced to leads. It is a \textbf{mutual information and decision system}. Both parties should gain enough accurate evidence to decide whether working together is sensible.
